\section{Certified hierarchical elimination}\label{sec:hierarchy}

\subsection{Retained interfaces and the adjacent-interface certificate}

Choose disjoint retained sets and let $I$ be their complement:
\begin{equation}\label{eq:retained}
R=B\mathbin{\dot\cup}R_1\mathbin{\dot\cup}\cdots\mathbin{\dot\cup}R_K,\qquad I=V\setminus R.
\end{equation}
Let $I_1,\ldots,I_J$ be the connected components of $G[I]$, with retained
boundaries $\partial I_j=N_G(I_j)\cap R$. The groups $R_k$ are indexed
consecutively in the retained order. Stage indices can have gaps after pruning
or coarsening, so they cannot replace these compact indices when testing
adjacency.

\begin{definition}[adjacent-interface certificate]\label{def:certificate}
The retained partition is \emph{certified} if
\begin{enumerate}[label=(\roman*),itemsep=1pt]
\item every edge of $G[R]$ with both endpoints outside $B$ joins vertices in the
  same group or in two consecutive groups; and
\item for every interior component, $\partial I_j\setminus B\subseteq R_k\cup R_{k+1}$
  for some $k$. A component may also touch a single group or none.
\end{enumerate}
\end{definition}

Given adjacency lists, the certificate can be checked in $O(m+|\mathcal E|)$
time: find the components of $G[I]$, scan their incident edges while recording
the smallest and largest compact group index met, and check the retained-to-retained
edges in the same scan. No table indexed by components and retained rows is
needed. With the ordering $I_1,\ldots,I_J,R$, the operator has the block-arrow
form
\begin{equation}\label{eq:arrow}
H_\delta=\begin{bmatrix}
A_1&&&G_1\\&\ddots&&\vdots\\&&A_J&G_J\\G_1^\top&\cdots&G_J^\top&D
\end{bmatrix},
\end{equation}
in which $G_j$ has nonzero columns only on $\partial I_j$.

\begin{theorem}[exact chain condensation]\label{thm:chain}
Suppose $H_\delta\succ0$ and \cref{def:certificate} holds. Then
\begin{equation}\label{eq:reduced}
S=D-\sum_{j=1}^{J}G_j^\top A_j^{-1}G_j\succ0
\end{equation}
is block tridiagonal on $R_1,\ldots,R_K$, with arbitrary coupling to $B$. For
every right-hand side $f$, solving $H_\delta x=f$ is equivalent to
\begin{equation}\label{eq:condense}
\widehat f_R=f_R-\sum_jG_j^\top A_j^{-1}f_j,\qquad
Sx_R=\widehat f_R,\qquad
x_j=A_j^{-1}(f_j-G_jx_R).
\end{equation}
\end{theorem}

\begin{proof}
Each $A_j$ is a principal submatrix of an SPD matrix and is therefore
invertible; block elimination gives~\eqref{eq:reduced} and~\eqref{eq:condense}.
The update from component $j$ is supported on $\partial I_j\times\partial I_j$.
By condition~(ii), it can fill only one retained group, two adjacent groups and
the border. Condition~(i) gives the same restriction for $D$. Hence no two
nonadjacent groups become coupled. Positive definiteness follows from the
Schur-complement property of SPD matrices.
\end{proof}

\Cref{eq:reduced} is an algebraic description. The implementation forms neither
$A_j^{-1}$ nor $A_j^{-1}G_j$; partial elimination of a local augmented matrix
produces the retained-front update directly (\cref{sec:local-factor}).
\Cref{fig:hierarchy} illustrates the structure.

\begin{figure}[htbp]
\centering
\begin{tikzpicture}[x=1cm,y=1cm,font=\small]
  % (a) graph view
  \foreach \x/\lab in {0/1,2.45/2,4.9/3}{
    \draw[fill=papergray!15,draw=papergray] (\x,0) rectangle ++(1.7,0.9);
    \node at (\x+0.85,0.45) {$I_{\lab}$};
  }
  \foreach \x/\lab in {1.95/1,4.4/2}{
    \draw[fill=paperblue!35,draw=paperblue] (\x-0.2,0) rectangle ++(0.4,0.9);
    \node[below] at (\x,0) {$R_{\lab}$};
  }
  \node at (6.95,0.45) {$\cdots$};
  \draw[fill=paperorange!25,draw=paperorange] (0,1.65) rectangle (7.3,2.15);
  \node at (3.65,1.9) {global border $B$};
  \foreach \x in {0.85,3.3,5.75,1.95,4.4}{
    \draw[densely dashed,papergray] (\x,0.9) -- (\x,1.65);
  }
  \draw[decorate,decoration={brace,mirror,amplitude=4pt},papergray]
    (0,-0.45) -- (6.6,-0.45) node[midway,below=4pt,align=center,font=\footnotesize]
    {interiors: eliminated independently, in parallel};
  \node[font=\footnotesize] at (3.65,2.55) {(a) retained interfaces and interior components};
  % (b) reduced matrix
  \begin{scope}[shift={(8.4,-0.6)}]
    \def\n{4}\def\w{0.72}
    \draw[papergray] (0,0) rectangle (3.5,3.5);
    \foreach \k in {0,...,3}{
      \fill[paperblue!45] ({\k*\w},{3.5-(\k+1)*\w}) rectangle ++(\w,\w);
    }
    \foreach \k in {0,...,2}{
      \fill[paperblue!18] ({(\k+1)*\w},{3.5-(\k+1)*\w}) rectangle ++(\w,\w);
      \fill[paperblue!18] ({\k*\w},{3.5-(\k+2)*\w}) rectangle ++(\w,\w);
    }
    \fill[paperorange!30] (4*\w,{3.5-4*\w}) rectangle (3.5,3.5);
    \fill[paperorange!30] (0,0) rectangle (4*\w,{3.5-4*\w});
    \fill[paperorange!60] (4*\w,0) rectangle (3.5,{3.5-4*\w});
    \draw[papergray] (0,0) rectangle (3.5,3.5);
    \node at (1.44,2.06) {$T$};
    \node at (3.21,1.8) {$Q$};
    \node at (3.21,0.33) {$C_B$};
    \node[font=\footnotesize,align=center] at (1.75,3.95) {(b) reduced matrix $S$};
  \end{scope}
\end{tikzpicture}
\caption{The certified hierarchy. (a) Removing the retained groups $R_k$ and the
border $B$ leaves interior components $I_j$ that touch at most two consecutive
groups. (b) Eliminating the interiors exactly yields a block-tridiagonal chain
$T$ on $R_1,\ldots,R_K$ coupled to a small dense border block $C_B$.}
\label{fig:hierarchy}
\end{figure}

\subsection{Separator pruning and temporal coarsening}

Let $\SR_R(H)$ denote the Schur complement of $H$ onto $R$, obtained by
eliminating $V\setminus R$.

\begin{proposition}[associativity of Schur complements]\label{prop:assoc}
If $R'\subseteq R$ and $H\succ0$, then $\SR_{R'}\bigl(\SR_R(H)\bigr)=\SR_{R'}(H)$.
\end{proposition}

\begin{proof}
For fixed $x_R$,
\begin{equation}\label{eq:minform}
x_R^\top\SR_R(H)\,x_R=\min_{x_{V\setminus R}}x^\top Hx.
\end{equation}
Minimizing further over $x_{R\setminus R'}$ is the same as minimizing the
original form over $x_{V\setminus R'}$. The quadratic forms coincide, and the
minimizers are unique because $H\succ0$.
\end{proof}

\Cref{prop:assoc} justifies two simplifications. \emph{Carry pruning} absorbs
separator rows that terminate in a one-sided branch into a neighboring interior,
and \emph{temporal coarsening} removes selected cuts so that several stages
form one interior. In both cases the new partition must again satisfy
\cref{def:certificate} before a chain solver is used. The heuristic that proposes
candidates is not trusted; only the certificate is. Reducing $|R|$ does not
necessarily reduce total work, because larger interiors can create larger
local fronts (\cref{sec:select}).

\subsection{Existence of a compatible hierarchy}

The certificate is not only an a posteriori test: under the modeling assumptions
a compatible hierarchy can be constructed directly.

\begin{theorem}[spaced-cut construction]\label{thm:spaced}
Suppose \cref{ass:local,ass:support} hold. Choose an integer $h\ge\ell$, divide
the stages into consecutive slabs of at most $h$ stages, place cuts at
$h,2h,\ldots$ strictly before $N$, and retain the rows $Q_t$ of
\cref{thm:monomial} at each cut. Then the retained sets are pairwise disjoint
and, after empty groups are removed, satisfy \cref{def:certificate}. Every
interior component lies in one slab. If each cut shares at most $s$ variables
and $P=\binom{s+2r}{2r}$, then
\begin{equation}\label{eq:spaced}
p\le P,\qquad b_j\le2P+g,\qquad n_j\le hq,\qquad K\le N/h.
\end{equation}
\end{theorem}

\begin{proof}
A row crossing the cuts at $t$ and $t+h$ has $a_i\le t$ and $b_i\ge t+h+1$, so
its span is at least $h+2>\ell$, which contradicts \cref{ass:support}; hence no
non-border row is retained twice. A non-retained row crosses no cut, so all of
its resources lie in one slab, and rows in different slabs share no resource.
Every interior component therefore lies in one slab, and since all of its rows
first appear in that slab, it has at most $hq$ vertices.

A retained neighbor of a component shares a resource inside the component's slab
and crosses its own cut. If that cut were not a boundary cut of the slab, the
row would also cross the boundary cut lying in between, and would be retained
twice. Hence a component touches only the two boundary groups of its slab and
$B$. Finally, two retained rows at nonconsecutive cuts cannot share a resource:
at a cut strictly between them, the earlier row has all of its resources on the
left and the later row all of its resources on the right, since otherwise one of
them would cross two cuts. Empty groups have no incident edges and can be
removed. The bounds follow from \cref{thm:monomial}, the slab argument and the
number of cuts.
\end{proof}

The construction proves existence and correctness. It does not minimize fronts,
and the solver usually retains smaller cuts and prunes more aggressively, using
the certificate to validate each choice.

\subsection{The rank-one correction}\label{sec:rankone}

Let $v=H_\delta^{-1}u$. The Sherman--Morrison identity gives
\begin{equation}\label{eq:sm}
\Hfull^{-1}f=H_\delta^{-1}f-v\,\frac{u^\top H_\delta^{-1}f}{1+u^\top v}.
\end{equation}
Since $H_\delta\succ0$, the denominator is at least one in exact arithmetic.
Preparing $v$ costs one extra base solve per factorization, and applying the
correction costs $O(m)$ per right-hand side. The vector $v$, the denominator, the
local factors and the reduced factors must all refer to the same matrix and
shift, and floating-point acceptance checks the complete operator including the
rank-one term (\cref{sec:acceptance}).
